\Level 2 {Efficiency and scaling}

\begin{frame}{Speedup}
  \begin{itemize}
  \item Single processor time $T_1$, on $p$ processors $T_p$
  \item speedup is $S_p=T_1/T_p$, $S_P\leq p$
  \item efficiency is $E_p=S_p/p$, $0< E_p\leq 1$
  \end{itemize}
Many caveats
\begin{itemize}
\item Is $T_1$ based on the same algorithm? The parallel code?
\item Sometimes superlinear speedup.
\item Can the problem be run on a single processor?
\item Can the problem be evenly divided?
\end{itemize}
\end{frame}

\begin{frame}{Limits on speedup/efficiency}
  \begin{itemize}
  \item $F_s$ sequential fraction, $F_p$ parallelizable fraction
  \item $F_s+F_p=1$
  \item $T_1 = (F_s+F_p)T_1 = F_sT_1 + F_pT_1 $
  \item Amdahl's law: $T_p = F_sT_1 + F_pT_1/p $
  \item $P\rightarrow\infty$: $T_P\downarrow T_1F_s$
  \item Speedup is limited by $S_P\leq 1/F_s$,
    efficiency is a decreasing function $E\sim 1/P$.
  \item loglog plot: straigth line with slope~$-1$
  \end{itemize}  
\end{frame}

\begin{frame}{Scaling}
  \begin{itemize}
  \item Amdahl's law: strong scaling\\
    same problem over increasing processors
  \item Often more realistic: weak scaling\\
    increase problem size with number of processors,\\
    for instance keeping memory constant
  \item Weak scaling: $E_p>c$
  \item example (below): dense linear algebra
  \end{itemize}
\end{frame}

